\documentclass[10pt,a4paper]{amsart}
\usepackage[margin=19mm]{geometry}
\usepackage{amsmath,amssymb,mathtools}
\usepackage[scr=rsfs,cal=euler]{mathalfa}
\usepackage{xfrac}
\usepackage[unicode,hidelinks,bookmarks=false]{hyperref}
\newcommand{\ie}{i.e.\ }
\newcommand{\cf}{cf.\ }
\newcommand{\resp}{respectively\ }
\newcommand{\Subsection}{\subsection}
\providecommand{\nobreakdash}{\nobreak\hbox{-}}
\setlength{\emergencystretch}{2em}
\multlinegap0pt
\usepackage{amssymb}
\usepackage{dsfont}
\usepackage{enumerate}
\newtheorem{lemma}{Lemma}
\newtheorem{prop}{Proposition}
\newtheorem{thm}{Theorem}
\newcommand{\N}{\mathbb{N}}
\newcommand{\V}{\mathcal{V}}
\newcommand{\W}{\mathcal{W}}
\newcommand{\id}{\text{id}}
\newcommand{\vr}{\nu}
\title{Garside shadows and biautomatic structures in Coxeter groups}
\author{Fabricio Dos Santos}
\date{Source: arXiv:2505.21718v2 (CC BY 4.0)}
\begin{document}
\maketitle

\noindent\emph{Source and licence.} Garside shadows and biautomatic structures in Coxeter groups. Version arXiv:2505.21718v2 (CC BY 4.0), distributed by arXiv under the Creative Commons Attribution 4.0 International licence, http://creativecommons.org/licenses/by/4.0/, as recorded in the arXiv licence metadata for this exact version and shown on the abstract page https://arxiv.org/abs/2505.21718v2. Full licence notice: \texttt{licenses/arxiv-2505.21718-CC-BY-4.0.txt}.\par\smallskip

\noindent\textit{Editorial scope.} Counted unit: the article's prop propFTPb (source0.tex lines 375--377). The article's complete original proof of it is delivered with it (source0.tex lines 379--409, one \texttt{proof} environment of 31 lines). No mathematics is written, repaired or re-derived: the statement and the proof are the article's own lines, in the article's own notation, and no retained line is substituted or re-flowed.  Retained uncounted as the source-local context the proof needs: lines 138--140; lines 154--156; lines 256--262; lines 347--349; lines 369--371.  One declared adaptation: exactly 1 ancillary strings inside the retained spans are omitted (image-inclusion commands and, where the source repeats a label, the repeated label definitions) and no other line differs from the source; the figure environments, with their placement, captions and labels, are kept, and the package records the omitted strings in fidelity receipts bound to the affected bodies.  No retained line contains a citation, so the wrapper carries no bibliography and no bibliography entry.  The retained lines contain the article's own diagram code, which the wrapper typesets with the standard tikz packages; no image file is delivered and the package declares no asset dependency.  Editorial formatting only, in addition: an AMS \texttt{amsart} preamble that restates the article's own declarations and macros used by the retained lines, namely the environments \texttt{lemma}, \texttt{prop}, \texttt{thm} on separate counters, together with the standard packages those lines need.  The retained environments print in this document as Lemma 1, Proposition 1, Theorem 1 in order of appearance, whereas the article prints the same items with the numbering of its own full document.  The complete native source of the article as distributed in the arXiv e-print for this exact version is bundled unchanged as \texttt{sources/178883/source0.tex}.
\begin{lemma} \label{refinement}
    Let $B$ and $B'$ be Garside shadows such that $B \subseteq B'$, and let $x,y \in W$. If $\pi_{B'}(x) = \pi_{B'}(y)$, then $\pi_B(x) = \pi_B(y)$. 
\end{lemma}

\begin{lemma} \label{lem-mShi}
    If $g \in W$ and $m \in \N$ then $\pi_{L_m}(g)$ is the gate of the $m$-Shi part containing $g$. Therefore, the partition $\mathscr{P}_{L_m}$ is equal to $\mathscr{S}_m$ and the partition $\mathscr{P}_L$ is equal to $\mathscr{S}$.
\end{lemma}

\begin{lemma} \label{constM}
Let $B$ be a finite Garside shadow. Then there exists a constant $M = M(B)$ such that:
\begin{enumerate}
    \item \label{contM-1} $B \subseteq L_M$;
    \item $d(\vr_B(g), g) \leq M$ for every $g \in W$.
\end{enumerate}
\end{lemma}

\begin{prop} \label{PropFTPa}
    If $B$ is finite then $\V_B$ satisfies the first fellow traveller property with $C = 2M$, where $M$ is the constant from Lemma \ref{constM}.
\end{prop}

\begin{thm} \label{Parallel wall}
    Let $W$ be a Coxeter group and $m \in \N$. There exists a constant $Q_m = Q(W,m)$ such that for each $g \in W$ and wall $\W$ at distance greater than $Q_m$ from $g$ in $X^1$, there are $m$ walls separating $g$ from $\W$.
\end{thm}

\begin{prop} \label{propFTPb}
    If $B$ is finite then $\V_B$ satisfies the second fellow traveller property with $C' = 4M(M+Q_M)+2Q_M$, where $M$ is the constant from Lemma \ref{constM} and $Q_M$ is the constant from Theorem \ref{Parallel wall}.
\end{prop}

\begin{proof}
    Let $g \in W$ and $s \in S$. We will prove this by induction on $\ell(g)$. Assume without loss of generality that $\ell(sg)>\ell(g)$. For $g = \id$ there is nothing to prove, so suppose $g \neq \id$. Let $v,v' \in \V_B$ be words representing $g$ and $sg$ respectively. 
    
    Suppose first that $\pi_B(g^{-1}) = \pi_B((sg)^{-1})$. Then,
    \begin{align*}
        s\vr_B(g) = sg\pi_B(g^{-1}) = sg\pi_B((sg)^{-1}) = \vr_B(sg)
    \end{align*}
    so $v'(\ell(\vr_B(sg)))$ and $sv(\ell(\vr_B(g)))$ represent the same element of $W$. Hence, the second fellow traveller property is satisfied by induction for $i < \ell(\vr_B(sg))$ and by Lemma \ref{constM} for $i \geq \ell(\vr_B(sg))$.

    Now, assume that $\pi_B(g^{-1}) \neq \pi_B((sg)^{-1})$. Consider the wall $\W_s$ separating $\id$ from $s$. Since $\vr_B(g) \preceq g$ we have that $s \preceq s\vr_B(g) \preceq sg$, so $s\vr_B(g)$ lies on the same side of $\W_s$ as $s$. If this is also the case for $\vr_B(sg)$, then $s \preceq \vr_B(sg)$ since we can reflect through $\W_s$ the part of the geodesic from $\id$ to $\vr_B(sg)$ in $X^1$ lying on the same side of $\W_s$ as $\id$. We thus have a geodesic $\gamma$ from $\id$ to $sg$ in $X^1$ passing through $s$ and $\vr_B(sg)$, but then translating $\gamma$ by $(sg)^{-1}$ we find that $\pi_B((sg)^{-1}) \preceq g^{-1} \preceq (sg)^{-1}$, which implies that $\pi_B(g^{-1}) = \pi_B((sg)^{-1})$ by definition of $\pi_B$, giving a contradiction. Hence, $\vr_B(sg)$ lies on the same side of $\W_s$ as $\id$. See Figure \ref{f:ftpfig} for an illustration.

    \begin{figure}[h]
    \begin{center}
    
    \end{center}
    \caption{Proof of Proposition \ref{propFTPb}.}
    \label{f:ftpfig}
    \end{figure}
    
    Finally, we claim that the wall $\W_s$ is $M$-close to $sg$. Assume for a contradiction that it is not. Then, $(sg)^{-1}\W_s$ is not $M$-elementary, and since it separates $g^{-1}$ from $(sg)^{-1}$, we must have that $\pi_{L_M}(g^{-1}) = \pi_{L_M}((sg)^{-1})$ as a consequence of Lemma \ref{lem-mShi}. But since $B \subseteq L_M$ by Lemma \ref{constM}, we have by Lemma \ref{refinement} that $\pi_B(g^{-1}) = \pi_B((sg)^{-1})$ which is a contradiction. Hence, $\W_s$ is $M$-close to $sg$. 

    Let $1 \leq i \leq \ell(v)$ and $\overline{v(i)} = h \in W$ so that $s \preceq sh \preceq sg$. Then, $\W_s$ is also $M$-close to $sh$ since any wall separating $\W_s$ from $sh$ also separates $\W_s$ from $sg$ (otherwise it intersects the geodesic $sv$ twice). Hence, by Theorem \ref{Parallel wall} we have that $d(sh, h) \leq 2Q_M$ and in particular $d(s\vr_B(g), \vr_B(g)) \leq 2Q_M$. Therefore,
    \begin{align*}
        d(sg, \vr_B(g)) \leq d(sg, s\vr_B(g)) + d(s\vr_B(g), \vr_B(g)) \leq M + 2Q_M
    \end{align*}
    and thus $d(\vr_B(g), \vr_B(sg)) \leq 2(M+Q_M)$. Hence, by applying Proposition \ref{PropFTPa} repeatedly we have for $i < \ell(\vr_B(sg))$ that $d(\overline{v(i)}, \overline{v'(i)}) \leq 4M(M+Q_M)$ and thus
    \begin{align*}
        d(s\overline{v(i)}, \overline{v'(i)}) \leq 4M(M+Q_M) + 2Q_M.
    \end{align*}
    Finally, since $d(s\vr_B(g), \vr_B(sg)) \leq 2M$, we have that $d(s\overline{v(i)}, \overline{v'(i)}) \leq 2M$ for all $i \geq \ell(\vr_B(sg))$.
\end{proof}

\end{document}
